\subsection{REFLECTIONS}

Recall that from now on, unless stated otherwise, the field K is assumed to be of characteristic different from 2. A reflection in V is a pseudo-reflection s such that \(s^{2}=1\) ; if s is a reflection, we denote by \(V_{s}^{+}\) the kernel of s-1 and by \(V_{s}^{-}\) that of \(s+1\) .

PROPOSITION 2. Let \(s\) be an endomorphism of V.

\begin{enumerate}
\def\labelenumi{(\roman{enumi})}
\item
  If \(s\) is a reflection, \(V\) is the direct sum of the hyperplane \(V_s^+\) and the line \(V_s^-\) .
\item
  Conversely, assume that V is the direct sum of a hyperplane H and a line D such that \(s(x) = x\) and \(s(y) = -y\) for \(x \in H\) and \(y \in D\) . Then s is a reflection and \(H = V_{s}^{+}\) , \(D = V_{s}^{-}\) . Finally, D is the image of 1 - s.
\item
  If \(s\) is a reflection, \(V_{s}^{+}\) is a hyperplane. If \(x\) belongs to \(V_{s}^{+} \cap V_{s}^{-}\) , then \(x = s(x) = -x\) , so \(x = 0\) since \(K\) is of characteristic \(\neq 2\) . Finally, for \(x\) in \(V\) , the vector \(x' = s(x) + x\) (resp. \(x'' = s(x) - x\) ) belongs to \(V_{s}^{+}\) (resp. \(V_{s}^{-}\) ) since \(s^2 = 1\) , and \(2x = x' - x''\) . Thus \(V\) is the direct sum of \(V_{s}^{+}\) and \(V_{s}^{-}\) , and \(V_{s}^{-}\) is necessarily of dimension 1 since \(V_{s}^{+}\) is a hyperplane.
\item
  Under the stated hypotheses, every element of V can be written uniquely in the form \(v = x + y\) with \(x \in H\) and \(y \in D\) and we have \(s(v) = x - y\) ; assertion (ii) follows immediately from this.
\end{enumerate}

COROLLARY. If V is finite dimensional, every reflection is of determinant -1.

Let \(s\) be a reflection in V. Prop. 2, (i) shows that there exists a basis \((e_1, \ldots, e_n)\) of V such that \(s(e_1) = e_1, \ldots, s(e_{n-1}) = e_{n-1}\) and \(s(e_n) = -e_n\) , hence \(\det(s) = -1\) .

PROPOSITION 3. Let \(s\) be a reflection in \(V\) .

\begin{enumerate}
\def\labelenumi{(\roman{enumi})}
\item
  A subspace \(V'\) of \(V\) is stable under \(s\) if and only if \(V_s^- \subset V'\) or \(V' \subset V_s^+\) .
\item
  An endomorphism \(u\) of \(V\) commutes with \(s\) if and only if \(V_s^+\) and \(V_s^-\) are stable under \(u\) .
\item
  If \(V' \subset V_{s}^{+}\) , then \(s(x) = x\) for all x in \(V'\) , so \(s(V') \subset V'\) . Assume that \(V_{s}^{-} \subset V'\) ; then, for any x in \(V'\) , \(s(x) - x \in V_{s}^{-} \subset V'\) , hence \(s(x) \in V'\) ; thus \(s(V') \subset V'\) . Conversely, assume that \(s(V') \subset V'\) ; if \(V' \not\subset V_{s}^{+}\) , there exists x in \(V'\) with \(s(x) \neq x\) ; the non-zero vector \(a = s(x) - x\) belongs to the line \(V_{s}^{-}\) , and hence generates this space; since \(a \in V'\) , we have \(V_{s}^{-} \subset V'\) .
\item
  Assume first that u commutes with s. If x is a vector such that \(s(x) = \varepsilon.x\) (where \(\varepsilon = \pm1\) ), then \(s(u(x)) = u(s(x)) = \varepsilon.u(x)\) , so \(V_{s}^{+}\) and \(V_{s}^{-}\) are stable under u. Conversely, assume that \(V_{s}^{+}\) and \(V_{s}^{-}\) are stable under u; it is clear that us - su vanishes on \(V_{s}^{+}\) and on \(V_{s}^{-}\) , and since V is the direct sum of \(V_{s}^{+}\) and \(V_{s}^{-}\) , we have us - su = 0.
\end{enumerate}

COROLLARY. Two distinct reflections \(s\) and \(u\) commute if and only if \(\mathrm{V}_s^- \subset \mathrm{V}_u^+\) and \(\mathrm{V}_u^- \subset \mathrm{V}_s^+\) .

If \(V_{s}^{-} \subset V_{u}^{+}\) and \(V_{u}^{-} \subset V_{s}^{+}\) , Prop. 3, (i) shows that \(V_{u}^{+}\) and \(V_{u}^{-}\) are stable under \(s\) , hence \(su = us\) by Prop. 3, (ii).

Conversely, if \(su = us\) , the subspace \(\mathrm{V}_{s}^{-}\) is stable under \(u\) by (ii); by (i), there are two possible cases:

\begin{enumerate}
\def\labelenumi{\alph{enumi})}
\item
  We have \(V_{u}^{-} \subset V_{s}^{-}\) : since they are both of dimension 1, these spaces are therefore equal, hence \(V_{s}^{-} \not\subset V_{u}^{+}\) ; since \(V_{u}^{+}\) is stable under s, we have \(V_{u}^{+} \subset V_{s}^{+}\) and so these two hyperplanes are equal. But then s = u, contrary to our assumptions.
\item
  We have \(V_{s}^{-} \subset V_{u}^{+}\) : the image of 1 - s is thus contained in the kernel of 1 - u, so \((1 - u).(1 - s) = 0\) . Since s and u commute, we have \((1 - s).(1 - u) = 0\) , in other words \(V_{u}^{-} \subset V_{s}^{+}\) .
\end{enumerate}

Remark. Let \(a \neq 0\) be in \(V\) and let \(a^*\) be a non-zero linear form on \(V\) : it follows from formula (1) that

\[
s _ {a, a ^ {*}} ^ {2} (x) = x + \left(\langle a, a ^ {*} \rangle - 2\right) \langle x, a ^ {*} \rangle . a
\]

and hence that \(s_{a,a^*}\) is a reflection if and only if \(\langle a, a^* \rangle = 2\) . In this case, we have \(s_{a,a^*}(a) = -a\) .

